\documentclass[12pt,a4paper]{article}  % Déclare la classe du document.
% Il existe 5 classes sous LaTeX : article, book, report, letter et slides
% Les options de classe sont entre crochets et permettent de faire des choix d'ordre général :
% - définir la taille de base des caractères avec 10pt, 11pt, 12pt, les commandes d'agrandissement 
% ou de réduction de la tailles des caractères ( \small \large ) se feront alors par rapport à cette base
% - définir la taille du papier:  a4paper,  a5paper, b5paper, executivepaper, legalpaper ou letterpaper
% Utiliser a4paper dès que la papier utilisé est de ce format c'est à dire ... tout le temps ;-)
% - utiliser des options de mise en page : 
%       ->   landscape passe en mode paysage pour l'ensemble du doccument
%       ->   onecolumn  option par défaut, le texte sera sur une seule colonne 
%       ->   twocolumn   pour un doccument sur deux colonnes, des réglages sont possibles (cf doc & net)
%       ->   oneside toutes les pages seront traités identiquement, par défaut avec la classe article
%       ->   twoside  mise en page différentes pour les pages pairs et impairs par défaut avec book
%       ->   openright et openany  pour gérer le commencement des chapitres dans la classe book
%       ->   titlepage et notitlepage indique si une nouvelle page doit être commencée après le titre du document.

% \usepackage permet de déclarer un module qui sera pris en compte dans la suite. 
% Les modules permettent d'étendre les fonctionnalité de LaTeX

%  --------------------------------------------------------------------------------------------------------
%  ------------------------------------------  Modules usuels  --------------------------------------
%  --------------------------------------------------------------------------------------------------------

\usepackage{amsmath} 	% Module indispensable qui renforce les capacité mathématique de LaTeX
\usepackage{amssymb}	% Met à disposition tout un tas de symboles, mathématiques ou non.
\usepackage{listings}         % pour le code informatique
\usepackage{emp}	         % The EMP package allows to encapsulate METAPOST files in LATEX sources.

\usepackage{numprint} % mise en forme des nombres

\usepackage{pifont}

\usepackage{gensymb}

%  --------------------------------------------------------------------------------------------------------
% ----------------------  packages pour fontes et caractères spéciaux  ---------------------
%  --------------------------------------------------------------------------------------------------------

\usepackage{eurosym}    % permet d'utiliser la typographie officielle du symbole euro
\usepackage{amsfonts}    % permet d'obtenir des caractères spéciaux pour les maths
\usepackage{amstext}       % pour la commande \text
\usepackage{amsfonts}	% Pour les lettres à double barre représentant les ensembles de nombres avec l'instruction \mathbb{N}
\usepackage{mathrsfs}	% Pour les lettres calligraphiés majuscules avec \mathscr{}

% \usepackage{amsthm}  % pour utiliser l'environement theorem sans numérotation
% \usepackage{hanging}  % permet l'indentation à g de un ou plusieurs paragraphe

\usepackage{pstricks,pst-all,pstricks-add,pst-tree,pst-node,pst-3dplot}

%  --------------------------------------------------------------------------------------------------------
%  --------------------------  Francisation et gestion des caractères sous mac -------------------
%  --------------------------------------------------------------------------------------------------------

% \usepackage{fontspec}
% \usepackage[french]{babel}% gestion du français
% \usepackage{textcomp}% caractères additionnels
% \usepackage{mathtools,amssymb,amsthm}% packages de l'AMS + mathtools
% \usepackage{mathpazo}
% \usepackage[scaled=0.95]{helvet}
% \usepackage{courier}

\usepackage[francais]{babel}   
% permet d'utiliser les règles typographiques françaises par exemple concernant les césures. 

%   !!!  ATTENTION   !!!    bien lire les lignes suivantes :
% sous Windows il peut être nécessaire d'activer la ligne suivante :
% \usepackage[latin1]{inputenc}   % (permet d'utiliser les caractères ISO 8859-1)

\usepackage[utf8]{inputenc} % unicode
\usepackage[OT1,T1]{fontenc} % permet d'utiliser directement les symboles accentués dans le texte sans avoir recours à un codage du type \'e, mais celui-ci reste possible

%  --------------------------------------------------------------------------------------------------------
% ------------------------------------Packages  de  mise en page -----------------------------------
%  --------------------------------------------------------------------------------------------------------

\usepackage{fancyhdr}	% en-têtes et pieds de pages
\usepackage{amsfonts}

 \usepackage{multicol} % permet l'utilisation du multicolonnage dans le texte. La syntaxe est \begin{multicols}{nb de colonnes}    ......  \end{multicols}
 
\usepackage{calc}            % permet de faire un minimum de calcul sur les arguments des commandes et macros (utilisé dans la def de l'environnement preuve )
\usepackage{epic}            % extension et simplification de  l'environnement picture
\usepackage{graphicx}    % permet d'inclure des graphique externe dans un document
\usepackage{variations}   % permet d'obtenir de joli tableau de variations
\usepackage{enumerate}  %permet de faire des listes personnalisées
\usepackage{arydshln}     %permet d'avoir des lignes en pointillés

% Par exemple :
% \begin{enumerate}[argument] où argument vaut :
% A : \item produira A. B.....
% a : \item produira a. b.
% I : \item produira I, II, IIi
% i : \item produira i, ii, iii..
% 1: \item produira 1, 2, 3
% EX i : \item produira EX i., EX ii,
% {A}-I \item produira A-I, A-II,...
% \end{enumerate}
\usepackage{framed}  % Crée 3 nouveaux environnements, qui se comportent comme des minipage de largeur \linewidth, 
% mais qui peuvent, de plus, se scinder entre plusieurs pages.
% framed : avec un cadre autour;
%  shaded : avec un fond coloré (il faut définir la couleur shadecolor);
%  leftbar : avec une barre le long du côté gauche.
\usepackage{multido}  % package pour utiliser une boucle itérative
% Sa syntaxe est la suivante : 
% \multido{variables}{nbiteration}{code}
% Le code sera ainsi répété nbiteration fois. Les déclarations de variables sont séparées par des virgules. Un déclaration prend la forme :
%    variable = valeurinitiale + increment  Exemple :
% \multido{\i=0+1}{21}{instruction à répéter } 


%  --------------------------------------  MACROS USUELLES  ----------------------------------------

\newcommand{\R}{\mathbb{R}} 
\newcommand{\N}{\mathbb{N}} 
\newcommand{\D}{\mathbb{D}} 
\newcommand{\Z}{\mathbb{Z}} 
\newcommand{\Q}{\mathbb{Q}} 
\newcommand{\C}{\mathbb{C}} 

%  -------------------------------------  MACROS SPECIFIQUES ------------------------------------

% Raised Rule Command:
%  Arg 1 (Optional) - How high to raise the rule
%  Arg 2            - Thickness of the rule
\newcommand{\raisedrule}[2][0em]{\leaders\hbox{\rule[#1]{1pt}{#2}}\hfill}


\newcounter{numexo}

\setcounter{numexo}{1}

\newcommand{\exercice}[1]{\vskip 0.6cm\noindent\rule[+2pt]{1cm}{0.5pt}\ {\large\bf \ Exercice \thenumexo {  #1 }}\raisedrule[+2pt]{0.5pt}\null\addtocounter{numexo}{1}\par\vskip 0,4cm}

% Utilisation de \exercice : 
% Pour avoir le titre avec son numéro faire \exercice{}
% Pour avoir en plus un titre faire par exemple \exercice{: les trois carrés}
\newcommand{\lignes}[1]{
\setlength{\unitlength}{1mm}
\vskip 0,4cm\noindent
\multido{\i=0+1}{#1}{%
     \noindent
     \begin{picture}(150,5)
     \linethickness{0,005mm}
     \put(-5,0){\dashline{1}(175,5)(5,5)}
     \end{picture}\vskip 0,001cm
     }
\vskip 0cm
}

% permet de mettre un rappel des consignes au début du devoir : le contenu peut varier selon si le devoir est à rendre sur copie double ou directement sur le sujet (désactiver ou activer des lignes)
\newcommand{\presentation}{
\begin{center}
\fbox{
\begin{minipage}{0.80\linewidth}
Consignes de présentation : 
 \vskip 0,1cm
  \begin{itemize}
  \item rendre le sujet avec la copie en indiquant nom, prénom et classe,
%  \item utiliser une ou plusieurs copies doubles,
  \item tracer un cartouche d'au moins 10 carreaux sur la première page du devoir,
  \item bien numéroter les exercices et les questions sur la copie, 
  \item sauter une ligne entre deux questions, écrire lisiblement et sans ratures.
%  \item souligner ou encadrer les réponses aux questions.
%  \item faire les éventuelles figures géométriques sur une feuille à part. 

\end{itemize}
\end{minipage}
}
\end{center}		

}


% ----------------------------------------- FORMAT DE PAGE  -----------------------------------------

\usepackage[left = 1.3cm,
	 	      right = 1.3cm,
		      top = 1.7cm,
		      headsep = 8mm,
		      bottom = 1.7cm,
		      headheight = 14.2 pt
		      ]{geometry}

\setlength{\parindent}{0mm}


\everymath{\displaystyle} % permet d'avoir tout les math en displaystyle
\renewcommand{\thesection}{\Roman{section}}  % permet de numéroter les sections en chiffres romains

\author{}
\title{}
\date{}
\pagestyle{fancy}

%\lhead{621}   %  Indiquer la classe ici  
\lhead{1ère spécialité maths}   %  Indiquer la classe ici  

%\rhead{DS 5}		% Indiquer la date pour rendre le devoir ici
\rhead{DS n$^{\circ}$02 - 17/11/2021}		% Devoir Sujet B

\fancyfoot[C]{- \thepage\ -} % bas de page : centre par exemple:  - \thepage\ -
\fancyfoot[L]{} % bas de page : gauche
\fancyfoot[R]{} % bas de page : droite

\lstset{language=Python}  % langage utilisé dans les parties code informatique

\lstset{frame = single
framerule = 1pt
showspaces = false,
showstringspaces = false,
numbers=left, numberstyle=\tiny}

% ----------------------------------------- correction  -----------------------------------------
\usepackage[most,skins,breakable]{tcolorbox}
\usepackage{bm}

\newcommand{\corr}[1]{\fcolorbox{red}{yellow!30}{\bm{#1}}}

\newtcolorbox{corrBox}{enhanced,boxrule=1pt,nobeforeafter,boxsep=0pt,left=\fboxsep,right=\fboxsep,colback=yellow!30,colframe=red!75!black}

\newcommand{\va}[1]{\lvert \ #1 \ \rvert}
\newcommand{\vaf}[1]{\left\lvert \ #1\  \right\rvert}

\usepackage[column=X]{cellspace}
\setlength{\cellspacebottomlimit}{4pt}
\setlength{\cellspacetoplimit}{4pt}
\setlength{\parindent}{2pt}
\setlength{\tabcolsep}{0.5cm}

\usepackage{enumitem}
\newlist{enumMet}{enumerate}{2}
\setlist[enumMet,1]{label=\arabic*), font=\bf\color{blue!75!black}}
\setlist[enumMet,2]{label=\alph*), font=\bf\color{blue!75!black}}

\begin{document}

\begin{center}

\fbox{\LARGE{\rule{0.5cm}{0pt}\rule[-0.25cm]{0pt}{0.9cm} DS n$^{\circ}$02 - Trigonométrie et équations du 2nd degré	\rule{0.5cm}{0pt}}}
\end{center}

%\begin{center}\Large\textbf{Sujet A}\end{center}
\begin{center}\textbf{2h - Barème indicatif (sur 20)}\end{center}
%\begin{center}\textbf{Les élèves avec un tier-temps ne traitent pas les questions avec  le symbole \ding{46}}\end{center}
\begin{center}\textbf{Calculatrice autorisée.\\Les résultats devront être justifiés (à l'aide de calcul ou de propriétés).}\end{center}

\vspace{-.3cm} 



% \exercice{ - valeur absolue - ( 4 points )}
% %environ 15-20 min
% Exprimer les nombres suivants sans le symbole valeur absolue.
% \begin{multicols}{2}
%     \begin{enumerate}
%         \item $\va{5-9}$
%         \item $\vaf{2-\dfrac{13}{5}}$
%         \item $\va{3-2\sqrt{2}}$
%         \item $\vaf{\dfrac{4-2\pi}{\pi-5}}$
%     \end{enumerate}
% \end{multicols}

% \exercice{ - valeur absolue  - ( 7 points )}
% $x$ est un réel tel que $-4\le x\le 2$. Déterminer un encadrement des expressions suivantes (attention les nombres des 2 bornes sont de signes différents).
% \begin{multicols}{2}
%     \begin{enumerate}
%         \item $\va{x}$ ;
%         \item $2\va{x}+1$ ;
%         \item $3\va{x}-2$ ;
%         \item $5-4\va{x}$;
%     \end{enumerate}
% \end{multicols}

% \exercice{ - résolution d'équations - ( 6 points )}
% Résoudre dans $\R$  les équations suivantes.
% \begin{multicols}{2}
%     \begin{enumerate}
%         \item $(2x-5)^2=64$ ;
%         \item $(-x+1)^2+2=11$ ;
%         \item $3(x-1)^2-12=63$ ;
%         \item $-2(5x+7)^2-35=0$;
%     \end{enumerate}
% \end{multicols}

% \exercice{ - Tableau de variations - ( 6 points )}
% Pour chaque fonction polynôme du second degré suivante, dresser son tableau de variations (les valeurs du tableau de variations doivent être justifiées).
% \begin{enumerate}
%     \item $f(x)=3(x+4)^2+5$
%     \item $g(x)=-5(x-5)(x+6)$
% \end{enumerate}

% \exercice {- points et paraboles - ( 5 points )}
% A partir des points fournis (qui appartiennent à une parabole représentative d'une fonction polynôme du deuxième degré) retrouver les 3 expressions de la fonction (forme développée, forme factorisée et forme canonique).\\
% Soit les points $A(-4;0)$, $B(3;0)$ et $D(0;6)$, 3 points appartenant à la courbe représentative de la fonction $f$.

\exercice {- enroulement autour du cercle trigo - ( 2 points )}

\lignes{2}

En utilisant la figure ci-dessous, donner les points du cercle qui correspondent aux réels suivants. On détaillera la décomposition du nombre.
\begin{multicols}{4}
\begin{enumerate}
    \item $1235\pi$
    \item $\dfrac{2435\pi}{6}$
    \item $\dfrac{-253\pi}{3}$
    \item $-\dfrac{123\pi}{4}$
\end{enumerate}
\end{multicols}
\begin{center}
\newrgbcolor{ududff}{0.30196078431372547 0.30196078431372547 1.}
\newrgbcolor{qqwuqq}{0. 0.39215686274509803 0.}
\psset{xunit=5.0cm,yunit=5.0cm,algebraic=true,dimen=middle,dotstyle=o,dotsize=5pt 0,linewidth=2.pt,arrowsize=3pt 2,arrowinset=0.25}
\begin{pspicture*}(-1.1,-1.1550625290940115)(1.1008598498816817,1.2224219662799085)
\psaxes[labelFontSize=\scriptstyle,xAxis=true,yAxis=true,Dx=1.,Dy=1.,ticksize=-2pt 0,subticks=2]{->}(0,0)(-1.98066,-1.1550625290940115)(1.1008598498816817,1.2224219662799085)
\pscircle[linewidth=2.pt](0.,0.){5.}
\psline[linewidth=2.pt,linestyle=dashed,dash=1pt 1pt,linecolor=red](0.,0.)(0.8660254037844387,0.5)
\psline[linewidth=2.pt,linestyle=dashed,dash=1pt 1pt,linecolor=red](0.,0.)(-0.8660254037844387,0.5)
\psline[linewidth=2.pt,linestyle=dashed,dash=1pt 1pt,linecolor=red](0.,0.)(-0.8660254037844386,-0.5)
\psline[linewidth=2.pt,linestyle=dashed,dash=1pt 1pt,linecolor=red](0.,0.)(0.8660254037844384,-0.5)
\psline[linewidth=2.pt,linestyle=dashed,dash=1pt 1pt,linecolor=blue](0.,0.)(0.7071067811865476,0.7071067811865475)
\psline[linewidth=2.pt,linestyle=dashed,dash=1pt 1pt,linecolor=blue](0.,0.)(-0.7071067811865475,0.7071067811865476)
\psline[linewidth=2.pt,linestyle=dashed,dash=1pt 1pt,linecolor=blue](0.,0.)(-0.7071067811865477,-0.7071067811865475)
\psline[linewidth=2.pt,linestyle=dashed,dash=1pt 1pt,linecolor=blue](0.,0.)(0.7071067811865474,-0.7071067811865477)
\psline[linewidth=2.pt,linestyle=dotted,linecolor=qqwuqq](0.,0.)(0.5,0.8660254037844386)
\psline[linewidth=2.pt,linestyle=dotted,linecolor=qqwuqq](0.,0.)(-0.5,0.8660254037844387)
\psline[linewidth=2.pt,linestyle=dotted,linecolor=qqwuqq](0.,0.)(-0.5,-0.8660254037844385)
\psline[linewidth=2.pt,linestyle=dotted,linecolor=qqwuqq](0.,0.)(0.5,-0.8660254037844386)
\begin{scriptsize}
\psdots[dotsize=4pt 0,dotstyle=*,linecolor=darkgray](0.,0.)
\rput[bl](0.00960106297279085,0.042848354149331562){\Large\black{$O$}}
\psdots[dotstyle=*,linecolor=ududff](1.,0.)
\rput[bl](1.027748218996967,0.03743267192752309){\Large\black{$I$}}
\psdots[dotstyle=*,linecolor=ududff](0.8660254037844387,0.5)
\rput[bl](0.8761092808657066,0.5265068007775197){\Large\black{$A$}}
\psdots[dotstyle=*,linecolor=ududff](0.7071067811865476,0.7071067811865475)
\rput[bl](0.7190546663726156,0.7350105663171459){\Large\black{$B$}}
\psdots[dotstyle=*,linecolor=ududff](0.5,0.8660254037844386)
\rput[bl](0.5105511264421327,0.8920653507495917){\Large\black{$C$}}
\psdots[dotstyle=*,linecolor=ududff](0.,1.)
\rput[bl](0.00960106297279085,1.0409966118493248){\Large\black{$J$}}
\psdots[dotstyle=*,linecolor=ududff](-0.5,0.8660254037844387)
\rput[bl](-0.5427979259339429,0.9028967151932086){\Large\black{$D$}}
\psdots[dotstyle=*,linecolor=ududff](-0.7071067811865475,0.7071067811865476)
\rput[bl](-0.7658857895025951,0.7377184074280502){\Large\black{$E$}}
\psdots[dotstyle=*,linecolor=ududff](-0.8660254037844387,0.5)
\rput[bl](-0.9256482421766014,0.5265068007775197){\Large\black{$F$}}
\psdots[dotstyle=*,linecolor=ududff](-1.,0.)
\rput[bl](-1.0994558275842002,0.020140513038427326){\Large\black{$I'$}}
\psdots[dotstyle=*,linecolor=ududff](-0.8660254037844386,-0.5)
\rput[bl](-1.0000187858048398,-0.5597649802551862){\Large\black{$G$}}
\psdots[dotstyle=*,linecolor=ududff](-0.7071067811865477,-0.7071067811865475)
\rput[bl](-0.850548494949918,-0.8072236335711421){\Large\black{$H$}}
\psdots[dotstyle=*,linecolor=ududff](-0.5,-0.8660254037844385)
\rput[bl](-0.620921440476689,-0.9859411468908217){\Large\black{$K$}}
\psdots[dotstyle=*,linecolor=ududff](0.,-1.)
\rput[bl](0.036679444781944466,-1.1094567257687464){\Large\black{$J'$}}
\psdots[dotstyle=*,linecolor=ududff](0.5,-0.8660254037844386)
\rput[bl](0.5376295082512864,-0.9205254646690133){\Large\black{$L$}}
\psdots[dotstyle=*,linecolor=ududff](0.7071067811865474,-0.7071067811865477)
\rput[bl](0.7488408863626846,-0.7445157924602379){\Large\black{$M$}}
\psdots[dotstyle=*,linecolor=ududff](0.8660254037844384,-0.5)
\rput[bl](0.9004798244939449,-0.51305963446988031){\Large\black{$N$}}
\end{scriptsize}
\end{pspicture*}
\end{center}

\exercice {- cosinus et sinus - ( 3 points )}

En utilisant le cercle trigonométrique ci-dessus, recopier et compléter le tableau suivant:\\
\begin{tabular}{ | Xc | Xc | Xc | Xc | Xc | Xc | Xc |}
    \hline
    $x$  & $\dfrac{2\pi}{3}$ & $-\dfrac{5\pi}{4}$ & $\dfrac{17\pi}{6}$ &  $-9\pi$ & $-\dfrac{11\pi}{2}$ & $\dfrac{7\pi}{4}$\\ \hline
    Point image  &   &   &   &   &   & \\ \hline
    $\cos(x)$  &   &   &   &   &   & \\ \hline
    $\sin(x)$ &   &   &   &   &   & \\
    \hline  
\end{tabular}

\ \\

Pensez à tourner la page

\newpage

\exercice {- même point image - ( 2 points )}
Pour chacun des nombres suivants, déterminer deux autres réels ayant le même point image lors de l'enroulement de la droite numérique
\begin{multicols}{2}
\begin{enumMet}
    \item $\dfrac{4\pi}{5}$
    \item $-\dfrac{5\pi}{6}$
\end{enumMet}
\end{multicols}

\exercice {- convertir en degrés - ( 1 points )}
Pour chacun des angles suivants, convertir l'angle en degré. On arrondira si nécessaire à 1 chiffre après la virgule
\begin{multicols}{2}
\begin{enumMet}
    \item $\dfrac{4\pi}{13}$ rad
    \item $\dfrac{2\pi}{5}$ rad
\end{enumMet}
\end{multicols}

\exercice {- convertir en radians - ( 1 points )}
Pour chacun des angles suivants, convertir l'angle en radian. On donnera la valeur exacte
\begin{multicols}{2}
\begin{enumMet}
    \item $72\degree$
    \item $355\degree$
\end{enumMet}
\end{multicols}

\exercice {- équations trigonométriques - ( 4 points )}

Résoudre les équations suivantes : 
\begin{enumMet}
    \item $\cos(x)=\dfrac{\sqrt{3}}{2}$ avec $x\in \left[0;\pi\right]$
    \item $\sin(x)=\dfrac{\sqrt{2}}{2}$ avec $x\in \left[\dfrac{3\pi}{2};\dfrac{9\pi}{2}\right]$
    \item $\sin(x)=\dfrac{\sqrt{3}}{2}$ avec $x\in \left[-3\pi;\pi\right]$
    \item $\sin(x)=\dfrac{17}{4}$ avec $x\in \left[-5\pi;5\pi\right]$
\end{enumMet}

\exercice{ - résolution d'équations - ( 8 points )}
Résoudre dans $\R$  les équations suivantes.
\begin{multicols}{2}
    \begin{enumMet}
        \item $3 x^{2}+5 x-2 = 0$
        \item $3 x^{2}-x+5=0$
        \item $-2 x^{2}+7 x+4=0$
        \item $-x^{2}+x-7 = 0$
        \item $16 x^{2}-40 x+25 = 0$
        \item $-3 x^{2}+5 x+8=0$
    \end{enumMet}
\end{multicols}




\end{document}

